\documentclass[main.tex]{subfiles}

% --- Document Body ---
\begin{document}
\clearpage
\thispagestyle{empty}
\mbox{}
\clearpage

\setcounter{page}{1}

\chapter{Introduction}\label{chap:introduction}
% * change from we to I 
% * add references 
\epigraph{If you will cling to Nature... to the little things that hardly anyone sees, and that can so unexpectedly become big and beyond measuring; if you have this love of inconsiderable things and seek quite simply, as one who serves, to win the confidence of what seems poor: then everything will become easier, more coherent and somehow more conciliatory for you.}{Rainer Maria Rilke}

\section{A brief review of quantum physics}
Quantum mechanics was developed to provide an accurate description of the physics of systems at the atomic and subatomic length scales. You might know of it from counterintuitive, and often misinterpreted, statements. A quantum system can exist in a superposition of multiple physical states. When observed, this superposition collapses, and the system is found in one of these multiple states according to some probability distribution. It is possible to know the momentum, or position of a quantum system, but not both at the same time. Matter exhibits both particle and wave-like characteristics. 
%In this thesis, I will not contribute much to justifying, or giving an intuition for, these statements. 

Dissertations concerned with quantum mechanics often first address its perception as eerie or disconcerting, and in my thesis I will follow this tradition. You might be familiar with Einstein's statement ``I, at any rate, am convinced that God does not throw dice", in response to the probabilistic nature underlying the Copenhagen interpretation of quantum mechanics. It is fair to say that quantum mechanics led to a fundamental change in our understanding of nature, and as such, it helps to look at previous groundbreaking theories and the reaction to them. 

Before Newton published his laws of mechanics in the book \emph{Principia} (1687), Cartesian physics---which explained gravity through the mechanical action of invisible vortices of matter---dictated the scientific worldview. In his book, Newton challenged this view, by presenting a proof that it was incompatible with Kepler's second and third laws of planetary motion~\cite{Newton1687Principia}.  Despite the widespread and near-immediate praise of \emph{Principia}, Newton's principles encountered, at the same time, resistance from peers such as Berkeley and Leibniz as well as from several continental philosophers~\cite{mcmullin2001impact}. 

An important point of contention was \emph{action at a distance}---gravitational force is exerted without contact (even through a vacuum), such as the force keeping the Earth in orbit around the Sun---which was perceived to possess an occult character~\cite{westfall1983never}. The notion that two objects could exert force on each other at a distance without contact was so counterintuitive at the time that his critics argued that a `cause' or agent is necessary for gravity and that Newton had not identified such a cause. Philosophers would discuss and critique the implications of, and assumptions behind \emph{Principia} for the centuries to come. In the decades after publication, however, Newton's principles were used to accurately model and predict natural phenomena---notably, the reappearance of Halley's comet in 1758---and \emph{Principia} garnered universal acclaim and adoption from the scientific community; open questions and all.

It is through a similar lens that we can view quantum mechanics. Developed over the $20^{\mathrm{th}}$ century, we are still grappling with fundamental aspects of the theory~\cite{hance2022does, bassi2023collapse}, and the consequences they have for our understanding of nature. In the meanwhile, it has proved to be singular in its ability to accurately model the physics of systems on the atomic scale. Precise calculations of energy levels of electrons in atoms~\cite{cowan2023theory}, predictions of novel states of matter (notably the Bose-Einstein condensate~\cite{dalfovo1999theory}) and the ability to measure time to a remarkably low uncertainty of $\sim 10^{-18}$~\cite{RevModPhys.87.637, PhysRevLett.133.023401}, were all enabled with quantum mechanical principles. It is in this spirit that my thesis takes on quantum mechanics---chiefly, in utilizing its principles to model the physics of atomic systems. 

Over the $20^{\mathrm{th}}$ century, scientists also moved beyond merely studying these systems to making them evolve and interact in a \emph{controlled} manner. The quantum phenomenon of tunneling---where an object passes through an energy barrier that is classically forbidden---is exploited for instance, to store information on modern solid state drives (SSDs)~\cite{zuolo2017}. Lasers, used in applications as wide ranging as surgery, welding and spectroscopy, exploit the phenomenon of stimulated emission~\cite{maiman1960, schawlow1958, hooker2010}. A new class of technologies, called quantum technologies, are being developed that require the ability for control on an even finer level~\cite{dowling2003, acin2018, degen2017, azuma2023, shannon2021, altman2021}. This includes quantum sensors~\cite{degen2017}, quantum networks~\cite{kozlowski2019towards} and quantum computers~\cite{ladd2010quantum}---the focus of this thesis.
\section{Physical systems and computation}\label{sec:phys_syst}
How does one then think of quantum computers? I believe it informative to start off with the `computer' aspect of these systems, by drawing analogies with the computers we currently possess---termed \emph{classical computers} in the quantum community. Each such computer has at its base an underlying physical system---as I will motivate, this seemingly trivial statement has profound consequences. I now present a vastly simplified description of classical computers meant to highlight the physics involved.  
% What problems are they meant to solve, how do they compare to present technology, and why are several countries (seemingly) competing with each other in a global quantum technology race~\cite{roberson2021talking}? 

The underlying system here is that of a \emph{transistor}---a semiconductor based device operating at room temperature. Such transistors have three terminals: the source, the drain and the gate. Information is stored in \emph{bits}, with each bit corresponding to a base 2 number. Consider for example the floating-gate transistors~\cite{hasler2016scaling, kahng1967floating}, utilized in portable flash drives and early SSDs. The two values (zero or one) of each bit are represented by two levels of electric charge stored in the gate~\cite{hasler2001overview}. The amount of charge controls the conductivity of a channel between the source and the drain---if the transistor conducts, one assigns the number $1$ to the transistor. If it does not conduct, one assigns the number $0$. Modern drives are often built of trillions of transistors, allowing them to store terabytes of information. 

How does one then carry out calculations, that is, manipulate the stored information? This is realized by passing bits to a \emph{logic gate}; a combination of transistors that perform operations on the input bits, and output the transformed bits. For instance, a $\mathrm{NOT}$ gate outputs $0(1)$ if the input was $1(0)$. Using such gates as building blocks, one can construct \emph{circuits}: in our computers, dedicated circuits exist for basic arithmetic operations. When one executes code in one's programming language of choice, this code is first translated to \emph{machine-level} operations involving executions of various circuits. In fact, all applications of computers, such as writing mails or playing video games, are (at the machine level) sequences of operations on circuits~\cite{patterson2016}. Bits can also be communicated efficiently from one location to the other---this technical achievement underlies the modern internet. 

While such transistors power most of our computers, alternative physical systems are also being studied to carry out computation and possess unique advantages for specific applications. One can use photons propagating in optical media to carry out computation---with the underlying system an optical transistor~\cite{kazanskiy2022optical, chen2013all}. Such devices exploit the low latency (afforded by the high propagation speed of light in optical materials~\cite{badloe2022}) and physical phenomena such as interference to carry out highly efficient matrix multiplications~\cite{zhou2022phot}. You might also recall that the modern internet is powered by optical fibers---cables that transmit bits stored in photons---laid out on ocean beds across the world, enabling rapid speed-of-light communication~\cite{amiri2019technical}. Often, such systems might enable entirely novel methods of computation. Consider the memristor, or the memory resistor~\cite{strukov2008missing}---a system with a resistance that changes continuously depending on the history of the current that flowed through it. These systems are being investigated for enabling brain-inspired, or neuromorphic computers~\cite{li2018review}. The key idea is that the memristors can model weighted connections between neurons, or \emph{synapses}, with neuroplasticity enabled by the variable resistance. 

As given away in the name, quantum computers utilize systems that act according to the laws of quantum physics---such as atoms cooled to near-zero Kelvin temperatures, in this thesis---for storing and manipulating information. In a quantum bit (or a \emph{qubit}), one maps the values $0, 1$ to two quantum states $\ket{0},\ket{1}$ of the atom, such as hyperfine levels or distinct electronic states~\cite{shannon2021}. As a direct consequence, while a bit either has a value of $0$ or $1$ (corresponding to separate physical states of the transistor), a qubit can exist in a \emph{superposition}, or a linear combination of $\ket{0}$ and $\ket{1}$---with its general quantum state written as 
\begin{equation}
    \ket{\psi} = \alpha\ket{0}+\beta\ket{1},
\end{equation}
with $\alpha,\beta$ as complex coefficients normalized such that $|\alpha|^2+|\beta|^2=1$. The reader might recall that $|\alpha|^2, |\beta|^2$ represent the probabilities that a given measurement collapses the state to either $\ket{0}$ or $\ket{1}$. Two qubits can similarly exist in a superposition of states $\ket{00},\ket{01},\ket{10},\ket{11}$. The state over $N$ qubits is hence described by $2^N$ coefficients---this exponential scaling means that quantum computers excel at representing high-dimensional data. As a concrete example, consider the exact calculations for determining the electronic structure of the molecule $\mathrm{C_3H_8}$, where large scale computations require $\sim 19$ terabytes (or $1.52\times 10^{14}$ bits) to store the physical state over 46 spin-orbitals~\cite{gao2024distributed}. On the other hand, a quantum computer requires 46 qubits to store such a state~\cite{mcardle2020quantum}, and the number of qubits (bits) needed scales linearly (exponentially) with the number of orbitals. Such a problem is closely connected to the problem of modelling interactions between drug molecules and proteins, and as a result quantum computers are being explored for applications in the pharmaceutical industry~\cite{blunt2022perspective}. 

Second, in the quantum framework we can act on each of these coefficients in parallel by applying a \emph{quantum gate} or a quantum circuit~\cite{nielsen2010quantum}. Hence, there is an intrinsic parallelism in the way that quantum computers operate. A lot of recent interest in quantum computing can be attributed to the fact that classes of problems that are classically intractable can be solved efficiently---exploiting this parallelism---on a quantum computer. Such problems include the simulation of complex physical systems for drug discovery, battery design and the optimization of large-scale systems relevant to logistics and energy management~\cite{bayerstadler2021}. Preliminary studies suggest that, due to the ability of quantum computers to represent and manipulate high-dimensional data, they might also be adept at tackling use cases in machine learning~\cite{alchieri2021}. We note, however, that current methods to encode and extract classical data (as required by machine learning) on a quantum computer often do not scale favorably with problem size~\cite{biamonte2017}.

\section{Current state of quantum technologies}
If by now you are sold on quantum computing, it is time for a disclaimer: quantum systems are notoriously challenging to build. To ensure that neutral atoms can be used as qubits, one must first cool them close to absolute zero, and trap individual atoms in optical traps or \emph{tweezers}~\cite{shannon2021}. To store quantum information in electronic states of atoms, one has to induce transitions between these states by applying a laser near-resonant with the transition frequency. Recall that quantum systems, upon interaction with the environment, might lose their quantum character, or \emph{coherence}, and can subsequently be described using classical physics. One must hence be able to preserve the coherence of a quantum state over the time scale required to carry out computations. In isolation, these are all impressive feats---as honored by the Nobel Prize in Physics of 1997~\cite{nobel1997} (cooling and trapping), of 2005~\cite{nobel2005} (quantum optics and coherence) and of 2018~\cite{nobel2018} (optical tweezers). The long-standing skepticism at the prospect of realizing large-scale quantum computers~\cite{dyakonov2020will, aaronson2004multilinear, kalai2016quantum} begins to make sense when one considers that many of these methods, and more, must be realized in one setup to move towards this goal. Then, that they have been implemented in current-generation atomic devices (admittedly at small scales and with varying amounts of errors) is, to the author, rather impressive. Notably, quantum algorithms---which require all of the aforementioned parts---have been demonstrated in experiment~\cite{Bluvstein_2023, reichardt2024, graham_multi-qubit_2022, PhysRevA.111.032611}.  

You might recall the classical computers from the middle of the $20^{\mathrm{th}}$ century, that took up half a room and were plagued by multiple vacuum tube failures that limited both the uptime and the reliability of the outcomes. Several breakthrough developments were needed to go from these primitive devices to the mature consumer devices of the 90s, and equally many to go from these devices to those we possess now, in our backpacks and our pockets. These developments were made across various fields---the theorists worked for instance on \emph{error correction}, or adding redundancy in stored information to systematically reduce loss~\cite{hamming1950error}. Experimentalists studied how to \emph{scale} up the devices---for instance to build processors with more transistors, directly increasing their processing power. For instance, ASML (near Eindhoven) builds extreme ultraviolet lithography machines~\cite{wu2007extreme} that print the intricate patterns of transistors onto silicon wafers, enabling high densities of 200 million transistors per square millimeter required for state-of-the-art chips. %Near Eindhoven lies ASML, the only company in the world that can build devices capable of manufacturing state-of-the-art chips, where more than 200 million transistors must be packed into a square millimeter. 



I posit that quantum computers are, in their trajectory of development, at the point classical computers were in mid $20^{\mathrm{th}}$ century. To name a few sources of error, atoms might be lost from tweezers. To realize two-qubit gates, one must excite atoms to high-lying (or \emph{Rydberg}) electronic states, which have a finite lifetime and can decay, causing information loss. Lasers can often be noisy, atoms might move around in a trap leading to Doppler shifts. The current generation of devices is often termed the noisy intermediate scale quantum (NISQ) generation~\cite{Preskill_2018}. As with classical computers, significant research effort is focused on quantum error correction~\cite{PRXQuantum.5.010333, PhysRevResearch.6.043253, pecorari2025high, reichardt2024, postema2025, pecorari2025high, stephens2014fault, brown2020parallelized, jahn2021holographic, panteleev2021quantum, PRXQuantum.5.010333, PhysRevA.108.062426} and mitigation~\cite{liao2024machine, quek2024, RevModPhys.95.045005,endo2018practical, strikis2021learning, o2023purification, bultrini2023unifying, PhysRevLett.131.210602, PhysRevLett.131.210601, PRXQuantum.3.010345, PhysRevLett.131.230601, PhysRevX.11.041036}---to carry out reliable computation in the presence of error mechanisms---as well as scaling up the devices by increasing the qubit counts.  

Neutral atoms are by no means the only competent architecture to build quantum computers. Other leading alternatives include superconducting~\cite{huang2020} and trapped ion platforms~\cite{bruzewicz2019trapped}, and many more platforms are being currently investigated~\cite{stern2013, zhang2019semiconductor, kloeffel2013, gershenfeld1998, childress2013diamond, o2007optical}. It is then worth asking, why use atoms?  One reason is that each atom (of a given species) is identical to the other atoms---in comparison to say superconducting qubits, which have to be manufactured. Second, atom based computers are scalable owing to their electrical neutrality---as demonstrated by the recent realization of a system with an industry-leading $6100$ qubits~\cite{manetsch2024}. There are many other compelling reasons and challenges associated with neutral atoms; which I will discuss throughout this thesis. 
\section{Scope of thesis}
As a theorist, my research can broadly be placed in the category of error characterization and mitigation. I am interested in understanding the various loss mechanisms associated with atoms in the NISQ era, and consequently in calculating loss rates associated with these mechanisms. When it comes to gates, an important metric we can calculate is the \emph{gate fidelity}, which measures the distance between the ideal operation and the operation realized on the physical device. This allows us to estimate the reliability and capabilities of NISQ devices, and is hence often reported by leading quantum research groups~\cite{PRXQuantum.6.020334, PhysRevLett.127.130501, graham_multi-qubit_2022, madjarov_high-fidelity_2020, PhysRevLett.121.123603, tsai2024, PhysRevLett.123.170503, PhysRevLett.121.240501, Evered_2023}. 

However, characterization alone is not sufficient to realize large-scale quantum computers---one must also find methods to mitigate the effect of the losses, that is, ways to improve gate fidelities. If you recall that gates are realized with the application of lasers inducing transitions between atomic states, several variables, or `knobs'---such as the laser intensity, laser frequency or the specific states one chooses to encode information---can play a profound role in ensuring reliable computation. Early papers in this field, for example, obtained the shortest possible---or \emph{time optimal}---laser pulses that implement a desired gate using numerical methods~\cite{PhysRevA.63.032308, PhysRevA.65.032301,d2002optimal,RevModPhys.76.1037}. Recently, time-optimal laser pulses were also devised to implement gates on neutral atoms~\cite{jandura_time-optimal_2022}. 

Put simply, most error mechanisms propagate in time---that is, the losses increase as time goes on. Hence, time optimal pulses can significantly improve the gate fidelities, and also lead to shorter algorithm runtimes. In this thesis, we propose and investigate gate schemes operating on time scales comparable to the time optimal pulses with additional properties, such as robustness to hardware noise, that are desirable for error mitigation. We further present strategies to reduce the total duration and decay-related losses at the circuit level, that are directly applicable to quantum algorithms. %In this thesis, I propose and investigate a gate scheme operating on a time scale comparable to . These pulses are additionally robust to deviations in control functions that arise in experiment. I also   

This thesis is structured into the following chapters: 
\begin{enumerate}
  \item \textit{\textbf{\textcolor{DeepPineGreen}{\nameref{sec:theory-chapter}}}}: We will introduce the theoretical concepts associated with atomic physics, loss mechanisms, and quantum algorithms that is crucial for understanding the following research chapters. We will also bridge the gap between the physical systems of atoms and the abstract concepts of quantum circuits and algorithms.

  \item \textit{\textbf{\textcolor{DeepPineGreen}{\nameref{sec:robust-chapter}}}}: In this chapter we model and analyse various error mechanisms associated with neutral atom quantum computers. Combining our model with numerical optimization, we then devise pulses to implement two-qubit $\mathrm{CZ}$ gates that retain a high fidelities in the presence of significant experimental noise. We also address the crucial question of selecting the Rydberg state and demonstrate that, using our error model, it is possible to select a suitable Rydberg state that minimizes losses, hence presenting another approach to improving fidelities.  
  \item\textit{\textbf{\textcolor{DeepPineGreen}{\nameref{sec:parametrized-chapter}}}}: Parametrized gates such as $k$-controlled phase or $\ckp$ gates are of importance for quantum compilation and variational algorithms on near-term devices. We address the challenge of implementing $\ckp$ gates on the neutral atom platform. We introduce a neural network-based optimization approach to generate continuous families of control pulses---a task challenging for traditional numerical methods. This strategy allows us to devise high-fidelity pulse families that operate near the quantum speed limit, offering significant improvements in speed and flexibility over both decomposition-based methods and prior fixed-angle pulse designs.
  \item \textit{\textbf{\textcolor{DeepPineGreen}{\nameref{chapter:varalg}}}}: The neutral atom platform possesses unique features such as long-range interactions and native multiqubit gates. Utilizing these features, we propose a novel algorithm implementation towards quantum chemistry, that is well suited to near-term quantum devices. We identify two use cases that neutral atoms are especially suited for, in comparison to other platforms---efficient implementations of variational ansatze, and a tensor network assisted state preparation routine. 
  \item \textit{\textbf{\textcolor{DeepPineGreen}{\nameref{chap:conclusion}}}}
\end{enumerate}

%As quantum computing is multidisciplinary, we will often move from atomic physics, to computer science and machine learning, to quantum chemistry.

%In the 60s, most users of the computers were research scientists, and they often cared about the hardware underlying the computer. If one wished to implement an algorithm on these computers, one needed to possess a crucial knowledge of  the underlying platform, the errors they are challenged by, and ways to mitigate or correct for such errors to enable reliable computation. %Indeed, a lot of progress arose from such considerations of error--- for instance, we started with vacuum tubes and moved to transistors once we realized they were more robust to errors (and importantly, more *scalable*---much smaller). The research field of *error correction* also arose from such considerations. The idea here of encoding redundancy. %The current user of such devices, by contrast, often does not care about the underlying semiconductors. They work with a high level framework. When sending messages, one does not fear the underlying hardware failing. One also does not have to %I would like to end this section by w. Quantum field theory--a landmark theory connecting special relativity and quantum mechanics--stated that the four fundamental forces are mediated by the exchange of bosonic particles, and hence provided a satisfying answer to the problem of `action at a distance'. At the same time, this problem would come back in a different form--that of quantum non-locality and the EPR paradox. 

\end{document}